\section{Meaning, evidence and continuation}\label{sec:interpretation}
\subsection{Capacity-sized does not mean whole-native-sized}\label{sec:amount}
The result links three events on one unchanged solution: P4-certified external access, a later receiver stock larger than a quarter of the joint capacity, and subsequent quantified net native recovery. Its denominators are specified in Table~\ref{tab:accounts}. In particular,
\begin{equation}\label{eq:initialfraction}
 \frac{C}{H_N(0)}=\frac23(1-\cos\phi)\longrightarrow0
 \qquad\text{as }\phi\downarrow0 .
\end{equation}
The ceiling $E_{\rm out}(t)\le C$ implies that both actual external stocks and the positive net recovered amounts in~\eqref{eq:recoveries} are at most $C$. Absolute stocks and recovered amounts can therefore tend to zero while the fractional guarantees remain uniform. The small phase interval is an essential restriction.

There is no claim of ceiling saturation, large whole-native transfer, an exact maximizing time, first return, nonlinear recurrence or period. The accounts are net changes rather than tagged-cohort transport or sustained monotonicity. The receiver is the original collocated conservative receiver, not an irreversible load, and it has no assigned material location. Trial-storage commensurability does not constitute physical calibration.

\subsection{Evidence and contribution status}\label{sec:evidence}
The accepted investigation and its actual GPT review are retained unchanged~\cite{investigation,review}. The investigation's fourteen exact-algebra groups were independently replayed during that internal review, with matching group results. A replay of those same fourteen groups is not twenty-eight separate findings. The written arguments handle regularity, periodic integration, zero-norm regularization, the scalar majorant, strictness and denominators; symbolic checks are not a machine-formal proof of these analytical steps.

The present work is manuscript preparation from accepted mathematics. GPT formulated and internally reviewed the investigation; Codex independently investigated it and prepared this manuscript under the research author's direction. The new PDF and proposed payload await GPT/author manuscript review. No external peer review is claimed, no scientific suite is rerun for typesetting, and retained source records are not described as fresh remote verification.

\subsection{Geometric distinctions and the next agreed direction}\label{sec:continuation}
The energy account does not reconstruct a spatial circuit. The programme retains the central starting/meeting region, opposed cones, outward spreading, curved around-and-inward return, stationary supports and possible transfer elsewhere as geometric directions. The adopted carrier circle and intended tight toroid remain distinct, as do the close original enclosing sphere and extended-line sphere. Finite host-notch N regions are not the three-upper/three-lower central G boundaries. The tentative 15--20\%, possibly more, remains author intuition rather than a coefficient or denominator.

The bounded receiving/recovery investigation is closed at its accepted scope. Broader phase/rate coverage, optimal transfer, saturation and physical/spatial calibration remain open questions, not defects to repair in this certificate. After manuscript review and separately authorized publication closeout, the agreed next scientific assessment is native $F$ stepping and exchange cadence, distinguishing the exact native baseline increment from exchange debit or credit. No native $F$ jump, cadence or physical clock enters this theorem.

Intended central-route correspondence and six-opening transport follow that direction. Gravity--magnetism stays linked across the programme; outer geometry, distinct golden-ratio relationships and further justified SRG integration remain later directions. Three-Way remains parked. Pi and H\'ok\'i p\'ok\'i brainstorming supply no premise here; the elementary inequalities in~\eqref{eq:pibounds} are only arithmetic tools. This discussion commissions no successor investigation.

